\documentclass[../../main.tex]{subfiles}
\graphicspath{{../../image/}{../../figure/}}
\begin{document}

\section{注记}\label{section:ch16-注记}

平衡律组 Cauchy 问题的局部与整体解，在\cref{theorem:ch16-1611}与\cref{theorem:ch16-1621}的假设下，由 Dafermos 与 Hsiao [1] 用随机选取法构造，随后由 Amadori 与 Guerra [1,3] 用波前追踪算法构造，最后由 Christoforou [1] 用消失黏性方法构造。特别地，Amadori 与 Guerra [1,3] 建立了解的唯一性，并把对角占优条件\eqref{align:ch16-s02-e04}改写成不依赖于 $DF$ 特定特征向量选择的形式。经由消失黏性的唯一性在 Christoforou [2] 中讨论。

关于消失黏性逼近的收敛速率，见 Christoforou 与 Trivisa [2]。正波的衰减速率在 Goatin 与 Gosse [1]、Christoforou 与 Trivisa [1,3] 中建立。带记忆系统的消失黏性途径在 Christoforou [3]、Chen 与 Christoforou [1] 中推进。关于初边值问题，见 Colombo 与 Guerra [3]。

关于\cref{theorem:ch16-1612}类似命题的详细证明（经由算子分裂），见 Christoforou [5]。或者，不使用算子分裂，可以把 Glimm 格式改造以适用于非齐次平衡律组，即在每一步求解一个广义 Riemann 问题（见原书第 9.11 节）。
在这方面，见 Hong 与 LeFloch [1]、Chou、Hong 与 Su [1]、Su、Hong 与 Chou [1]。

\cref{theorem:ch16-1613}取自 Tai-Ping Liu [14]。Amadori、Gosse 与 Guerra [1]、Seung-Yeal Ha [1] 通过建立 $L^1$ 稳定性改进了这一结果。波与源项之间共振的效应可见于 Tai-Ping Liu [18]、Cai Zhong Li 与 Tai-Ping Liu [1]、Pego [4]、Isaacson 与 Temple [4]、Klingenberg 与 Risebro [2]、Ha 与 Yang [1]、Lien [1]、Lan 与 Lin [1]、Asakura [3]、Hong 与 Temple [1]、Hong 与 Temple [1]。

Euler 方程的径向对称形式原书(7.1.29)（可能带阻尼、燃烧、自引力或 Euler--Poisson 方程的其他表现），提供了非齐次平衡律组的有趣例子。
大量相关问题已由 Gui-Qiang Chen [6]、Chen 与 Glimm [1,2]、Gui-Qiang Chen 与 Tiang-Hong Li [1]、Hsiao、Luo 与 Yang [1]、Tong Yang [1]、Wang 与 Wang [2,3]、Chen 与 Wagner [1]、Tsuge [1,2,3,4,5,6]、Ha、Huang 与 Lien [1] 研究。在上述多数工作中，为避免系统在 $r=0$ 处的奇异行为，必须把原点排除在解的定义域之外。见原书第 18.9 节。

阻尼重分布由 Dafermos [23,25,34] 引入。特别地，\cref{section:ch16-原点锥奇异熵l1稳定}中阐释的系统原书(5.5.52)的 $L^1$ 估计的推导细节，见 Dafermos [25]。
取自 Dafermos [36,43] 的\cref{section:ch16-部分耗散源l1稳定}中 $L^1$ 界的推导，本质上利用了 Ruggeri 与 Serre [1] 发明的 Liapunov 泛函。\cref{theorem:ch16-1665}的详细证明在 Dafermos [41]。关于周期解的\cref{theorem:ch16-1666}来自 Dafermos [42,44]。阻尼重分布在黏弹性与热传导中出现的带记忆平衡律组上的应用，见 Dafermos [38,39,40]。综述见 Dafermos [45]。

\end{document}
